\section{Praktikum}
\label{sec:praktikum-bab3}

\begin{praktikumbox}
\vspace{1.1em}
\textbf{Judul:} Konversi bilangan dan pemeriksaan data di memori.\\[0.3em]
\textbf{Sub-CPMK:} 1.1 \quad\textbullet\quad
\textbf{Alokasi waktu:} 100 menit \quad\textbullet\quad
\textbf{Bentuk:} individu
\end{praktikumbox}

\subsection{Tujuan Praktikum}

Setelah menyelesaikan praktikum ini, mahasiswa mampu mengonversi
bilangan antar basis, menuliskan nilai dalam beberapa basis pada
program assembly, serta memeriksa dan menafsirkan isi memori dengan
DEBUG dan emu8086.

\subsection{Alat dan Bahan}

\begin{itemize}
    \item Komputer dengan sistem operasi Windows.
    \item emu8086.
    \item DOSBox dengan TASM, TLINK, dan DEBUG.
    \item Berkas \file{bilangan.asm} dari \file{code/asm/}.
\end{itemize}

\subsection{Langkah Kerja}

\begin{enumerate}
    \item \textbf{Konversi manual.}
    \begin{enumerate}
        \item Salin tabel berikut ke lembar kerja.
        \item Lengkapi kolom biner, oktal, dan heksadesimal tanpa
              memakai kalkulator.
    \end{enumerate}

    \begin{center}
    \begin{tabular}{@{}cccc@{}}
        \toprule
        \textbf{Desimal} & \textbf{Biner} & \textbf{Oktal} &
        \textbf{Heksadesimal} \\
        \midrule
        37  & & & \\
        45  & & & \\
        100 & & & \\
        156 & & & \\
        200 & & & \\
        243 & & & \\
        255 & & & \\
        \bottomrule
    \end{tabular}
    \end{center}

    \item \textbf{Menulis program data.}
    \begin{enumerate}
        \item Buka \file{bilangan.asm} pada editor emu8086 atau editor
              teks.
        \item Perhatikan empat baris data: \code{desimal},
              \code{biner}, \code{heksa}, dan \code{karakter}.
        \item Jelaskan pada lembar kerja mengapa tiga baris pertama
              menghasilkan byte yang sama.
    \end{enumerate}

    \item \textbf{Perakitan dan pengaitan.}
    \begin{enumerate}
        \item Buka DOSBox, lalu masuk ke direktori kerja.
        \item Rakit dengan perintah \file{tasm bilangan.asm}.
        \item Link dengan perintah \file{tlink bilangan.obj}.
        \item Jalankan program dengan perintah \file{bilangan}.
    \end{enumerate}

    \item \textbf{Pemeriksaan memori dengan DEBUG.}
    \begin{enumerate}
        \item Jalankan \file{debug bilangan.exe}.
        \item Ketik \file{d ds:0} untuk menampilkan isi segmen data.
        \item Catat delapan byte pertama yang muncul.
        \item Tandai byte yang mewakili \code{desimal}, \code{biner},
              \code{heksa}, dan \code{karakter}.
    \end{enumerate}

    \item \textbf{Modifikasi data.}
    \begin{enumerate}
        \item Ubah nilai data menjadi:
              \code{desimal DB 37}, \code{biner DB 11011011b},
              \code{heksa DB 0DBh}, dan \code{karakter DB 'Z'}.
        \item Rakit ulang, lalu periksa kembali isi memori.
        \item Tafsirkan byte \hex{DB} sebagai bilangan tidak bertanda
              dan sebagai bilangan bertanda 8 bit.
    \end{enumerate}

    \item \textbf{Pengamatan dengan emu8086.}
    \begin{enumerate}
        \item Buka program yang sama pada emu8086.
        \item Jalankan dengan \file{Single Step} sampai instruksi
              \lstasm|MOV DS, AX| selesai.
        \item Buka jendela \file{Memory} pada alamat \code{DS:0000}.
        \item Bandingkan tampilan emu8086 dengan hasil DEBUG.
    \end{enumerate}
\end{enumerate}

\subsection{Hasil yang Dikumpulkan}

\begin{enumerate}
    \item Tabel konversi bilangan yang telah dilengkapi.
    \item Berkas \file{bilangan.asm} versi modifikasi.
    \item Tangkapan layar hasil \file{d ds:0} pada DEBUG.
    \item Tangkapan layar jendela memori emu8086.
    \item Analisis singkat: mengapa \hex{F3} dapat berarti 243 atau
          $-13$, dan mengapa \code{'A'} disimpan sebagai \hex{41}.
\end{enumerate}

\begin{asesmenbox}
\vspace{1.1em}
\textbf{Asesmen praktikum 1 (bobot 5\% dari komponen praktikum).}

\begin{tabular}{@{}p{6.2cm}cp{4.6cm}@{}}
    \toprule
    \textbf{Indikator} & \textbf{Bobot} & \textbf{Kriteria} \\
    \midrule
    Ketepatan konversi bilangan & 35\% &
    Seluruh baris tabel benar atau kesalahan dapat dijelaskan \\
    \addlinespace
    Program data berhasil dirakit dan dijalankan & 25\% &
    Berkas \file{.EXE} terbentuk dan program selesai tanpa \eng{hang} \\
    \addlinespace
    Ketepatan pembacaan isi memori & 25\% &
    Empat byte data dikenali dan ditafsirkan dengan benar \\
    \addlinespace
    Analisis bilangan bertanda dan ASCII & 15\% &
    Penjelasan \hex{F3} dan \code{'A'} tepat serta beralasan \\
    \bottomrule
\end{tabular}
\end{asesmenbox}
